\documentclass{article}
\usepackage{geometry}
\geometry{margin=1in}

\title{Editing the USER CODE: Explanation}
\begin{document}

\maketitle

\section*{Overview}

The USER CODE section in the generated X-CUBE-AI project allows for custom implementation of data acquisition, preprocessing, inference, and post-processing. The functions \texttt{acquire\_and\_process\_data} and \texttt{post\_process} are the key points for integrating your application’s logic with the generated AI model.

\section*{1. One Input and One Output Only}
The generated code from X-CUBE-AI defines the number of input and output buffers as \texttt{AI\_NETWORK\_IN\_NUM} and \texttt{AI\_NETWORK\_OUT\_NUM}, both set to \texttt{1}:
\begin{itemize}
    \item This means that the network can process only a single input and produce a single output at a time.
    \item The input and output are handled as single-dimensional arrays, simplifying the structure but limiting throughput.
    \item For this reason, we have to feed data sequentially, one value at a time, which aligns with the design of the predefined float inputs.
\end{itemize}

\section*{2. Predefined Float Indices Array}
Instead of using a dynamic input buffer, we define a static array of float values to represent the sentiment scores for our text examples:
\begin{verbatim}
const float predefined_inputs[] = {0.8f, 0.5f, 0.2f};
\end{verbatim}
\begin{itemize}
    \item Each value in the array corresponds to a sentiment score mapped in the Python preprocessing step.
    \item These values are fed into the neural network input buffer one by one in the \texttt{acquire\_and\_process\_data} function.
    \item This approach eliminates the need for complex string or integer tokenization on the embedded side.
\end{itemize}

\section*{3. Editing \texttt{int acquire\_and\_process\_data(ai\_i8* data[])}}
This function is responsible for loading the predefined input data into the model’s input buffer:
\begin{itemize}
    \item The \texttt{data} parameter is a pointer to the input buffer where data needs to be written.
    \item The function loops through the \texttt{predefined\_inputs[]} array and assigns each value to the buffer, one at a time.
    \item For example:
    \begin{verbatim}
    float* input_data = (float*)data[0];
    input_data[0] = predefined_inputs[current_index];
    \end{verbatim}
    \item After loading the input value, the \texttt{current\_index} is incremented, allowing the next value to be used in the following inference.
\end{itemize}

\section*{4. Editing \texttt{int post\_process(ai\_i8* data[])}}
This function handles the output generated by the model:
\begin{itemize}
    \item The output buffer contains a single float value, which represents the prediction for the given input.
    \item Example of extracting and printing the output:
    \begin{verbatim}
    float* output_data = (float*)data[0];
    printf("Prediction: %f\r\n", output_data[0]);
    \end{verbatim}
    \item The function can be extended to map this numerical output back to a text label, such as:
    \begin{verbatim}
    if (output_data[0] < 0.3) {
        printf("Prediction: Negative Sentiment\r\n");
    } else if (output_data[0] < 0.7) {
        printf("Prediction: Neutral Sentiment\r\n");
    } else {
        printf("Prediction: Positive Sentiment\r\n");
    }
    \end{verbatim}
\end{itemize}

\section*{5. Reviewing \texttt{app\_x-cube-ai.c}}
The \texttt{app\_x-cube-ai.c} file is the primary user entry point for managing inference. Key sections include:
\begin{itemize}
    \item **Initialization**: Functions to set up the model and memory.
    \item **Inference Loop**: Where \texttt{acquire\_and\_process\_data} is called, followed by inference execution, and then \texttt{post\_process}.
    \item **Utility Functions**: Additional helper functions for managing buffer allocations, error handling, and memory cleanup.
\end{itemize}
Editing this file allows you to customize how data flows through the network and how results are displayed or processed.

\section*{Conclusion}
The USER CODE section in the X-CUBE-AI project is where application-specific logic is implemented. By editing the \texttt{acquire\_and\_process\_data} and \texttt{post\_process} functions, you can manage inputs and outputs effectively, ensuring the model integrates seamlessly with your application.

\end{document}
